% Generated draft: requires mathematical and source review.
\documentclass[11pt]{article}
\usepackage[margin=25mm]{geometry}
\usepackage{fontspec}
\setmainfont{Latin Modern Roman}
\usepackage{amsmath,amssymb,mathtools,unicode-math}
\setmathfont{Latin Modern Math}
\usepackage{enumitem,needspace,xurl,hyperref}
\setlength{\parindent}{0pt}
\setlength{\parskip}{5pt}
\setlength{\emergencystretch}{3em}
\newcommand{\N}{\mathbb{N}}
\newcommand{\Z}{\mathbb{Z}}
\newcommand{\Q}{\mathbb{Q}}
\newcommand{\R}{\mathbb{R}}
\newcommand{\C}{\mathbb{C}}
\newcommand{\F}{\mathbb{F}}
\newcommand{\A}{\mathbb{A}}
\newcommand{\PP}{\mathbb{P}}
\newcommand{\E}{\mathbb{E}}
\newcommand{\eps}{\varepsilon}
\newcommand{\dd}{\,\mathrm{d}}
\begin{document}
\section*{UnsolvedMath ID: 11000303}
\textbf{Source problem code:} catalogue entry
\subsection{Definitions and mathematical statement}
Let $H_{\Q}=H_1(\Sigma_g;\Q)$ with its symplectic intersection pairing and let $U_{\Q}=\Lambda^3H_{\Q}/H_{\Q}$, where the usual embedding of $H_{\Q}$ is given by wedging with the symplectic class. Let $\mathcal M_g$ and $\mathcal M_{g,*}$ be the closed-surface and one-marked-point mapping class groups. Write $R^*$ for their rational cohomological tautological rings, in the customary convention generated by Miller--Morita--Mumford classes and, in the pointed case, the marked-point Euler/psi class. The source genus and class conventions need confirmation.

Let $\mathrm{Sp}$ denote the symplectic invariance in the source. Use the pairing to identify $H_{\Q}$ with its dual when needed; the extraction alternates between dual and primal notation. Define
\[
 A_*={\left(\Lambda^*((\Lambda^3H_{\Q})^*)/J_*\right)}^{\mathrm{Sp}},\qquad
 A={\left(\Lambda^*U_{\Q}^*/J\right)}^{\mathrm{Sp}},
\]
where $J_*$ is the homogeneous ideal generated by the source-designated quadratic modules $[1^2]_{\mathrm{torelli}}\oplus[2^2]$, and $J$ is the ideal generated by its designated $[2^2]$. The labels describe symplectic representations, but their precise copies/embeddings in the exterior squares are not given here. The subscript ``torelli'' must not be discarded or replaced by every copy with the same abstract representation type.

\textbf{Conjecture 2.1.} The stated natural cohomology homomorphisms induce isomorphisms
\[
 A_*\xrightarrow{\cong}R^*(\mathcal M_{g,*}),\qquad
 A\xrightarrow{\cong}R^*(\mathcal M_g).
\]
Moreover, $A_*$ and $A$ are Poincar\'e duality algebras of formal dimensions $2g-2$ and $2g-4$, respectively. Here formal dimension $D$ means top nonzero graded degree $D$, with top piece $\Q$ and perfect product pairings in complementary degrees; it is not the total vector-space dimension. Both the isomorphisms and duality conclusions are part of the conjecture. The source does not reproduce the construction of the natural maps or the relation-module embeddings, so these remain dependencies of a fully explicit presentation.

For comparison, the source recalls that for one Riemann surface $X$, $H^*(\operatorname{Jac}(X);\Q)$ has formal duality dimension $2g$ and
\[
 H^*(\operatorname{Jac}(X);\Q)/([1^2])\cong H^*(X;\Q),
\]
where $[1^2]=\ker(\Lambda^2H_{\Q}^*\to\Q)$ is the primitive part under contraction with the pairing, and the parentheses denote its generated ideal. It views the conjecture as a family analogue. It also writes $\Lambda^*U_{\Q}^*$ as the cohomology associated with the primitive third cohomology $PH^3(\operatorname{Jac})$, of exterior-algebra formal dimension $\binom{2g}{3}-2g$, and rewrites the unpointed quotient with the $[2^2]$ ideal. This is an exterior-algebra/associated-torus interpretation: the ordinary singular cohomology of a vector space $PH^3$ itself would not have that form. The omitted geometric meaning of that shorthand must be specified before treating it as a literal space.
\subsection{Short English statement}
The proposed symplectic-invariant quadratic quotients of exterior algebras should present the pointed and unpointed tautological cohomology rings, with perfect duality in top degrees $2g-2$ and $2g-4$. The precise relation modules and natural maps are missing from the excerpt.
\subsection{Sources}
\begin{itemize}
\item Local UnsolvedMath registry, permanent ID 11000303, source code catalogue entry; no separate proposer or proposal year supplied.
\item Farb, edited volume \emph{Problems on Mapping Class Groups and Related Topics}, Conjecture 2.1, PDF page 360, in source text bearing S. Morita's chapter attribution; \url{https://www.math.uchicago.edu/~farb/papers/mcgbook.pdf}. The registry records PDF-font-structure extraction and access on 2026-07-29; the inventory distinguishes 329 labeled headings by source page.
\item No additional precise primary citation or modern status reference is supplied; the research summary states that its literature information is truncated.
\end{itemize}
\end{document}
